\chapter{Метод Фурье и представление через функцию Грина}
\section{Устранение неоднородности границ}
Для варианта 11 решаем
\[
 u_t=a^2u_{xx}-bu+f,\quad u(0,t)=\mu(t),\quad -u_x(l,t)=\nu(t),
 \quad u(x,0)=\phi(x),\qquad a>0,\ b\ge0.
\]
Пусть $\mu,\nu$ достаточно регулярны (например, $C^1$ по времени),
а $\phi\in L^2(0,l)$ и источник допускает интегральное представление ниже.
Линейная функция
\[
 \boxed{w(x,t)=\mu(t)-x\nu(t)}
\]
удовлетворяет обеим неоднородным границам. При $u=w+v$ получаем
\[
 v_t=a^2v_{xx}-bv+\widetilde f,\quad v(0,t)=0,\quad v_x(l,t)=0,
 \quad v(x,0)=\widetilde\phi(x),
\]
\[
 \widetilde f=f-\mu'-b\mu+x(\nu'+b\nu),
 \qquad \widetilde\phi=\phi-\mu(0)+x\nu(0).
\]
Такой переход является индивидуальным аналогом замены из лекции 24.
\section{Спектральная задача}
Ищем $-X''=\lambda X$, $X(0)=0$, $X'(l)=0$.
Интегрирование по частям даёт
$\lambda\int_0^lX^2dx=\int_0^l(X')^2dx>0$
для ненулевой собственной функции. При $\lambda=k^2$
имеем $X=C\sin(kx)$ и $\cos(kl)=0$. Поэтому
\[
 k_n=\frac{(n+\tfrac12)\pi}{l},\quad \lambda_n=k_n^2,
 \quad X_n(x)=\sin(k_nx),\quad N_n=\frac l2,\quad n=0,1,\ldots.
\]
Здесь нет ни нулевой, ни отрицательной моды.
\section{Модальные уравнения}
Обозначим
\[
 \sigma_n=a^2k_n^2+b,\qquad
 c_n=\frac2l\int_0^l\widetilde\phi(\xi)\sin(k_n\xi)\,d\xi,
\quad F_n(t)=\frac2l\int_0^l\widetilde f(\xi,t)\sin(k_n\xi)\,d\xi.
\]
Разложение $v=\sum_nT_nX_n$ даёт
\[
 T_n'+\sigma_nT_n=F_n,\quad T_n(0)=c_n,
\qquad T_n(t)=c_ne^{-\sigma_nt}+\int_0^te^{-\sigma_n(t-\tau)}F_n(\tau)\,d\tau.
\]
Следовательно,
\[
 \boxed{u(x,t)=\mu(t)-x\nu(t)+\sum_{n=0}^{\infty}\sin(k_nx)
 \left[c_ne^{-\sigma_nt}+\int_0^te^{-\sigma_n(t-\tau)}F_n(\tau)\,d\tau\right].}
\]
Если $f_n=(2/l)\int_0^lf(\xi,t)\sin(k_n\xi)d\xi$, то
\[
 F_n=f_n-\frac{2(\mu'+b\mu)}{lk_n}
       +\frac{2(-1)^n(\nu'+b\nu)}{lk_n^2}.
\]
Использованы интегралы
$\int_0^l\sin(k_n\xi)d\xi=1/k_n$,
$\int_0^l\xi\sin(k_n\xi)d\xi=(-1)^n/k_n^2$.
\section{Функция Грина}
Введём тепловое ядро
\[
 \boxed{G(x,\xi,t)=\frac2l\sum_{n=0}^{\infty}
 \sin(k_nx)\sin(k_n\xi)e^{-(a^2k_n^2+b)t},\quad t>0.}
\]
Тогда окончательная формула имеет вид
\[
\boxed{\begin{aligned}
 u(x,t)={}&\mu(t)-x\nu(t)
 +\int_0^lG(x,\xi,t)[\phi(\xi)-\mu(0)+\xi\nu(0)]\,d\xi\\
 &+\int_0^t\int_0^lG(x,\xi,t-\tau)
 [f(\xi,\tau)-\mu'(\tau)-b\mu(\tau)
       +\xi(\nu'(\tau)+b\nu(\tau))]\,d\xi\,d\tau.
\end{aligned}}
\]
Начальные данные восстанавливаются в $L^2$, а при достаточной гладкости
и согласовании --- в классическом смысле. Формула учитывает произвольные
индивидуальные неоднородности; подставлять гармоники нужно лишь при
решении соответствующего практического пункта.
